% Daftar Singkatan dan Lambang (List of Abbreviations and Symbols)
% — Juknis p.11-12. Page numbers (first use) cannot be determined
% without an actual compile in this environment; left as TODO
% rather than guessed. No mathematical symbols (LAMBANG) beyond
% the statistical notation already defined inline in Chapter 3
% (e.g. p, N) are used densely enough in this thesis to warrant a
% separate symbol table — only the abbreviation table is populated.
\chapter*{Daftar Singkatan dan Lambang}
\addcontentsline{toc}{chapter}{Daftar Singkatan dan Lambang}

% Column widths kept within the 14cm text block (21cm A4 minus the
% 4cm/3cm margins): l (~1.5cm) + p{7cm} + p{3cm} = ~11.5cm, with
% headroom. A prior revision used p{9cm}+r (unwrapped), which
% overflowed the margin by ~83pt (~2.9cm) under the new 4cm/3cm
% geometry — found by an actual compile, fixed here.
\begin{center}
\begin{tabular}{lp{7cm}p{3cm}}
\toprule
\textbf{Singkatan} & \textbf{Nama} & \textbf{Halaman Pemakaian Pertama} \\
\midrule
FER & \textit{Facial Expression Recognition} & \TODO{p.} \\
CLAHE & \textit{Contrast Limited Adaptive Histogram Equalization} & \TODO{p.} \\
CVAT & \textit{Computer Vision Annotation Tool} & \TODO{p.} \\
LR & Logistic Regression & \TODO{p.} \\
RQ & Research Question / Pertanyaan Penelitian & \TODO{p.} \\
LAB & \textit{L*a*b* color space} & \TODO{p.} \\
\bottomrule
\end{tabular}
\end{center}
